% BÀI TOÁN HÌNH HỌC - BAI_HINH_25 (TRÍCH ĐỀ 25)
\begin{ex}
Cho đường tròn tâm O, bán kính R. Từ một điểm M ở ngoài đường tròn, kẻ hai tiếp tuyến MA và MB với đường tròn (A, B là các tiếp điểm). Gọi E là trung điểm của đoạn thẳng MB, C là giao điểm của AE và (O) ($C \neq A$), H là giao điểm của AB và OM. Gọi D là giao điểm của MC và (O) ($D \neq C$).
	\begin{enumerate}
		\item[a)] Chứng minh: Các điểm A, M, B, O cùng thuộc một đường tròn.
		\item[b)] Chứng minh: $BE^2 = EC \cdot EA$ và tam giác ABD là tam giác cân.
	\end{enumerate}
\loigiai{
{
	\begin{center}
		\includegraphics[scale=1.3]{hinh/25}
	\end{center}
}
\end{ex}
